% Original research note for integration into the parent report.
% This file is an input fragment, not a standalone preamble.
\section{An additional improvement from the actual Cauchy row domain}
\label{sec:annulus-modes}

The factorial Cauchy packet in the reviewed proposals uses the enlarged
row domain $z\in2\mathbb Z_{(2)}$ (up to a highly divisible extension-ring
perturbation). The actual reciprocal rows exclude a residue class.
Keeping this information gives a termwise stronger integral mode basis
without adding any source-tile families. The result is local: its worst
limiting slope over all valuation levels remains three.

Write $\nu_2(0)=+\infty$, and let
\[
 S_i=\{w\in\mathbb Z_{(2)}:0\le\nu_2(w)\le i\},\qquad i\ge0.
\]
Define $G(b)=b+\nu_2(b!)$ for $b\ge0$. The next construction is explicit;
no theorem about generalized factorials is needed.

\begin{lemma}[An explicit ordering of a punctured two-adic ball]
\label{lem:annulus-order}
For each $i\ge0$ there are distinct positive integers
$a^{(i)}_0,a^{(i)}_1,\ldots$ in $S_i$ and nondecreasing integers
$F_i(0),F_i(1),\ldots$ such that
\[
 F_i(b)=\min_{w\in S_i}\sum_{h<b}\nu_2(w-a^{(i)}_h),
 \qquad
 \sum_{h<b}\nu_2(a^{(i)}_b-a^{(i)}_h)=F_i(b).
\]
They can be constructed as follows. For $i=0$ take
$a^{(0)}_b=2b+1$ and $F_0(b)=G(b)$. For $i\ge1$, the sequence
$F_i$ is the sorted merge, with multiplicity, of the two strictly
increasing sequences
\[
 G(n),\qquad n+F_{i-1}(n),\qquad n=0,1,2,\ldots.
\]
When the next element comes from the first sequence output the point
$2n+1$; when it comes from the second output $2a^{(i-1)}_n$.
Resolve a tie in favor of the first sequence. For every $i,b$,
\[
 a^{(i)}_b\le2b+1,
 \qquad F_i(b)\ge\nu_2(b!).
\]
\end{lemma}

\begin{proof}
For the ordinary consecutive ordering $0,1,\ldots$, the minimum
valuation of a product of $n$ consecutive differences over
$\mathbb Z_{(2)}$ is $\nu_2(n!)$. This follows by counting, for every
$r\ge1$, the at least $\lfloor n/2^r\rfloor$ factors in the appropriate
residue class; equality holds at the next consecutive integer.
Scaling and translating gives the assertion for odd integers, including
rational two-integral arguments.

Split $S_i$ into its odd part $1+2\mathbb Z_{(2)}$ and its even part
$2S_{i-1}$. Differences between the two parts are units. If $n_o$
odd and $n_e$ even points have already been chosen, the least possible
valuation for the next odd point is $G(n_o)$, and that for the next even
point is $n_e+F_{i-1}(n_e)$. Choosing the smaller proves the asserted
minimum and the merge recursion. Both component sequences are strictly
increasing: $G$ increases by at least one, as does $n+F_{i-1}(n)$.

For the lower bound, use induction on $i$. The second component is
termwise at least $G$ once $F_{i-1}(n)\ge\nu_2(n!)$. The sorted merge
of two copies of $G$ is precisely $\nu_2(b!)$, because
$\nu_2((2n)!)=\nu_2((2n+1)!)=G(n)$. Thus the merge in the statement
is no smaller. The base case is immediate.

The same termwise comparison and the odd-first tie rule imply that
before an even point is selected, $n_o\ge n_e+1$. For an odd point,
$2n_o+1\le2b+1$. For an even point, induction bounds it by
$2(2n_e+1)\le2(n_o+n_e)=2b$. This proves the point-size bound and
also makes the construction polynomial in the number of requested
points and $i$.
\end{proof}

\begin{lemma}[Uniform valuation slope]
\label{lem:annulus-slope}
Put
\[
 R_i=1+\frac{3}{4^{i+1}-1}.
\]
Then, for all $i,b\ge0$,
\[
 R_i b-4\log_2(b+2)\le F_i(b)\le R_i b.
\]
\end{lemma}

\begin{proof}
The base case follows from
$F_0(b)=2b-s_2(b)$, where $s_2(b)$ is the binary digit sum.
For the induction step put $B=1+R_{i-1}$, so $2<B\le3$ and
\[
 R_i=\left(\frac12+\frac1B\right)^{-1}
     =\frac{2B}{B+2}.
\]
Substitution gives the displayed closed form for $R_i$.
The two component sequences are bounded above by $2n$ and $Bn$.
Their merge at index $b$ is at most $R_i b$: the number of terms in
these two arithmetic sequences no greater than $R_i b$ is at least
$b+1$, since the two unrounded ratios sum to the integer $b$.

For the lower bound let $x=F_i(b)$ and
$L_b=\log_2(b+2)$. Each of the first $b+1$ terms of the merge has a
component index at most $b$. The lower bounds for those components
are $2n-L_b$ and $Bn-4L_b$. Consequently
\[
 b+1\le \frac{x+L_b}{2}+1+\frac{x+4L_b}{B}+1,
\]
so
\[
 x\ge R_i b-R_i-\left(\frac{R_i}{2}
                      +\frac{4R_i}{B}\right)L_b.
\]
Here $R_i\le6/5$ and
$R_i/2+4R_i/B=(B+8)/(B+2)\le5/2$.
Since $L_b\ge1$, the last two losses together are at most
$4L_b$. This proves the uniform estimate.
\end{proof}

\begin{lemma}[Stability under an unramified extension perturbation]
\label{lem:annulus-extension}
Let $\mathcal O$ be the unramified coefficient ring of the contraction
argument, with valuation extending $\nu_2$. Suppose $R\ge2$,
$w_0\in S_i$, $w-w_0\in2^R\mathcal O$, and
$0\le b<2^{R-2}$. Then
\[
 \nu_2\left(\prod_{h<b}(w-a^{(i)}_h)\right)\ge F_i(b).
\]
\end{lemma}

\begin{proof}
The assertion is immediate for $b=0$. Otherwise the selected integers
are distinct and satisfy $1\le a^{(i)}_h\le2b-1<2^{R-1}$.
Choose the ordinary residue $c\in\{0,\ldots,2^R-1\}$ of $w_0$.
If $c$ is nonzero and is not a selected point, take $w'=c$.
It belongs to $S_i$, and all factor valuations are below $R$ and
unchanged by the perturbation. If $c=0$, then $\nu_2(w_0)\ge R$,
so $i\ge R$; take $w'=2^R\in S_i$. Again all differences to the
selected points have valuation below $R$.
Finally, if $c=a^{(i)}_h$ for one selected point, take
$w'=a^{(i)}_h+2^R\in S_i$. Its difference from that point has
valuation $R$, and its other difference valuations agree with those
of $w$; the corresponding difference for $w$ has valuation at least
$R$. In each case the product valuation for $w$ is at least that for
$w'\in S_i$, which is at least $F_i(b)$ by
Lemma~\ref{lem:annulus-order}. This argument uses only ordinary
residues of $w_0$; $w$ itself may have nonzero extension coordinates.
\end{proof}

\begin{theorem}[Annulus Cauchy modes]
\label{thm:annulus-modes}
At a contraction node of valuation $i$, retain the existing packet
layout and column tiles. Put
\[
 \alpha=2^{a+1},\quad T=t+\alpha\zeta,\quad
 R=a+1-i,\quad J=2^{a-i-1},\quad y=y_0+2^{i+1}j,
\]
where $0\le j<J$ and $\zeta\in\mathcal O$. For a row crossing
that tile, let
\[
 d_x=T-y_0+x,\quad z_x=\frac{2^{i+1}}{d_x},\quad
 A_x=\frac{\mathcal C(T)}{\mathcal C'(0)d_x},
 \qquad \xi_b=2a^{(i)}_b.
\]
Define
\[
 r_b(z)=\frac{b!\prod_{h<b}(z-\xi_h)}
                    {\prod_{h=0}^{b}(1-hz)},
 \qquad
 s_b(j)=\frac{(1-b\xi_b)\binom{j}{b}}
                    {\prod_{h=0}^{b}(1-\xi_h j)}.
\]
The exact finite identity
\[
 W_{xy}(T)=A_x\sum_{b=0}^{J-1}r_b(z_x)s_b(j)
\]
holds. All source functions $s_b(j)$ are two-integral; all displayed
new denominators are units. The target coefficient $A_xr_b(z_x)$ is
divisible by $2^{\omega_i(b)}$, where
\[
 \boxed{\ \omega_i(b)=R+b+F_i(b)+\nu_2(b!)\ }.
\]
In particular these weights are termwise no smaller than the previous
factorial weights with the same base contribution $R$, and
\[
 \omega_i(b)=R+\left(3+\frac{3}{4^{i+1}-1}\right)b
                 +O(\log(b+2)),
\]
with an absolute error constant. For $i=0$ the exact formula is
$\omega_0(b)=R+2b+2\nu_2(b!)$.
\end{theorem}

\begin{proof}
The crossing condition gives $e=\nu_2(t-y_0+x)\le i$. Since
$\alpha\zeta$ has valuation at least $a+1$, one has
$\nu_2(d_x)=e$, so $z_x$ is divisible by two and
$\nu_2(A_x)\ge a+1-e\ge R$, as in the original integrality proof.
For
\[
 w=z_x/2,\qquad w_0=2^i/(t-y_0+x),
\]
one has $w_0\in S_i$ and
\[
 \nu_2(w-w_0)\ge i+a+1-2e\ge a+1-i=R.
\]
Lemma~\ref{lem:annulus-extension} therefore yields
\[
 \nu_2\left(\prod_{h<b}(z_x-\xi_h)\right)
 \ge b+F_i(b)\qquad (b<J=2^{R-2}).
\]
The remaining denominator factors in $r_b$ are units, proving the
weight bound after multiplying by $b!$ and $A_x$.

For completeness the identity is a telescoping Cauchy elimination.
Let
\[
 U_b(z)=\prod_{h<b}\frac{z-\xi_h}{1-hz},\qquad
 V_b(j)=\prod_{h<b}\frac{j-h}{1-\xi_hj}.
\]
Then direct multiplication gives
\[
 \frac{U_b(z)V_b(j)-U_{b+1}(z)V_{b+1}(j)}{1-zj}
 =\frac{(1-b\xi_b)U_b(z)V_b(j)}{(1-bz)(1-\xi_bj)}
 =r_b(z)s_b(j).
\]
For $0\le j<J$ the factor $V_J(j)$ vanishes. Summation gives
$1/(1-zj)=\sum_{b<J}r_b(z)s_b(j)$, which becomes the kernel
identity after multiplication by $A_x$. Since every $\xi_b$ is even,
all source denominators and $1-b\xi_b$ are units. The binomial
coefficient is evaluated as an ordinary integer: no factorial is
divided in the residue ring. The comparison with the earlier weights
uses $F_i(b)\ge\nu_2(b!)$, and the slope follows from
Lemma~\ref{lem:annulus-slope} and
$\nu_2(b!)=b-s_2(b)$.
\end{proof}

\paragraph{Algorithmic consequence and its scope.}
Use one mode for each existing packet, tile, and index $b<J$, with
source linear form $\sum_j s_b(j)X_{y_0+2^{i+1}j}$. Its source and
target groups are unchanged. Hence the number of families, child
network sizes, entry-evaluation mechanism, and canonical-lift
invariant are unchanged. Truncate by the new exact weights.
The admissible tuple set is contained in that certified by the old
factorial weights, for every finite parameter choice. Precomputing
the ordering and evaluating the integral features costs only a
polynomial factor in the existing parameters. The uniform entropy
analysis may still use slope three, while a fixed valuation level
has the sharper slope
$\gamma_i=3+3/(4^{i+1}-1)$. These slopes approach three as $i$ grows.
Thus this result by itself is not a new logarithmic power in the
worst-case SAT exponent, nor does it prove a uniform leading constant
strictly better than the slope-three envelope over all levels.
A valuation-dependent taper may exploit the local improvements, but
that optimization is not assumed in the global theorem here.
